\subsection{The cubic branch reduces to rational codegree three or four}
\label{sec:unbounded-cubic-codegrees}
\begin{proposition}\label{prop:cubic-rational-codegree-bound}
Let $R$ be a self-dual rank-four fusion ring with
$R_{\mathbb Q}\cong\mathbb Q\times K$, where $K$ is a totally real cubic
field. Let $m$ be the formal codegree of the rational character. If $R$
admits a pivotal categorification, then $m=3$ or $m=4$.
More precisely, $m=1,5$ are impossible already at the ring level, while
$m=2$ or $m\geq6$ forces $D=\operatorname{FPdim}(R)<196$ and multiplicity
at most thirteen, and is excluded by the bounded classification.
\end{proposition}
\begin{proof}
The rational character has integer values, with first value one, so its
formal codegree $m$ is a positive integer. If $m=1$, its vector is
$(1,0,0,0)$, which cannot be orthogonal to the positive Frobenius--Perron
character. If $m=5$, its nontrivial entries have squared sum four, so the
vector is $(1,\pm2,0,0)$ up to permutation. Orthogonality to the positive
dimension vector would require $1\pm2d_i=0$, contrary to $d_i\geq1$.

By Proposition~\ref{prop:rank4-galois-exception}, a pivotal categorification
has a pseudo-unitary Galois conjugate or is already pseudo-unitary. If its
Frobenius--Perron character is rational, Corollary~\ref{cor:integral-selfdual-rank4}
already excludes the present algebra type. We may therefore assume that
it belongs to the cubic orbit. The formal codegrees are then
$m,D,D_2,D_3$, counted with their character multiplicities. The identity
for their reciprocals gives
\[
 D_2^{-1}+D_3^{-1}=1-\frac1m-\frac1D.
\]
Writing $w$ as in Lemma~\ref{lem:codegree-moment-bound}, Cauchy--Schwarz
therefore gives
\begin{align*}
 w&\geq\frac1{m^2}+
 \frac12\left(1-\frac1m-\frac1D\right)^2+\frac1{D^2}\\
 &=\left[\frac1{m^2}+\frac12\left(1-\frac1m\right)^2\right]
   -\frac{1-1/m}{D}+\frac3{2D^2}.
\end{align*}
For $m=2$, the bracket is $3/8$. For $m\geq6$, put $u=1/m\leq1/6$;
the function $u^2+(1-u)^2/2$ decreases on $[0,1/6]$ and equals $3/8$
at $u=1/6$. Thus in both cases
\[
 w\geq\frac38-\frac1D\geq\frac38-\frac2D.
\]
Corollary~\ref{cor:dimension-196} now gives
$D<196$, hence multiplicity at most thirteen. The bounded-classification
census has only one positive self-dual ring with rational algebra
$\mathbb Q\times\text{(cubic field)}$, that of
$\operatorname{PSU}(2)_7$, whose rational formal codegree is three.
It therefore excludes every case $m=2$ or $m\geq6$ under consideration.
Only $m=3,4$ remain for the unrestricted investigation.
\end{proof}
